\documentclass{article}

\usepackage{microtype}
\usepackage{graphicx}
\usepackage{subfigure}
\usepackage{booktabs}
\usepackage{hyperref}

\newcommand{\theHalgorithm}{\arabic{algorithm}}
\usepackage[accepted]{icml2024}

\usepackage{amsmath}
\usepackage{amssymb}
\usepackage{mathtools}
\usepackage{amsthm}

\usepackage[capitalize,noabbrev]{cleveref}

\theoremstyle{plain}
\newtheorem{theorem}{Theorem}[section]
\newtheorem{proposition}[theorem]{Proposition}
\newtheorem{lemma}[theorem]{Lemma}
\newtheorem{corollary}[theorem]{Corollary}
\theoremstyle{definition}
\newtheorem{definition}[theorem]{Definition}
\newtheorem{assumption}[theorem]{Assumption}
\theoremstyle{remark}
\newtheorem{remark}[theorem]{Remark}

\icmltitlerunning{Bilingual Sentiment Classification}

\begin{document}

\twocolumn[
\icmltitle{Ordinal Sentiment Classification of Bilingual Product Reviews}

\icmlsetsymbol{equal}{*}

\begin{icmlauthorlist}
\icmlauthor{Ismail Lataoui}{equal,eth}
\icmlauthor{Teammate Two}{equal,eth}
\icmlauthor{Teammate Three}{equal,eth}
\end{icmlauthorlist}

\icmlaffiliation{eth}{Department of Computer Science, ETH Zurich, Switzerland}

\icmlcorrespondingauthor{Ismail Lataoui}{ilataoui@ethz.ch}

\icmlkeywords{Sentiment Classification, NLP, Ordinal Regression}

\vskip 0.3in
]

\printAffiliationsAndNotice{\icmlEqualContribution}

\begin{abstract}
% TODO: write last. 4 sentences: problem, why hard, what we do, what we find.
We address ordinal sentiment classification of bilingual (German/English)
product reviews, where each review is mapped to a star rating in
$\{0,\dots,4\}$ and the evaluation metric is mean absolute error (MAE).
We compare three representations of increasing richness---sparse TF--IDF,
frozen multilingual sentence embeddings, and a fine-tuned multilingual
transformer---and show that adapting the encoder to the task is the key
driver of performance. We further introduce an MAE-optimal decoding rule
that exploits the ordinal structure of the metric at no training cost.
% TODO: final numbers
Our best model reaches a validation score of $0.91$, matching the strong
baseline.
\end{abstract}

\section{Introduction}
\label{sec:intro}
Product reviews are a primary signal for consumers and sellers alike, and
automatically recovering the star rating a reviewer assigned is a
long-standing sentiment-analysis task. We study a variant with three
properties that together make it non-trivial. First, the corpus is
\emph{bilingual}: reviews are split roughly evenly between German and
English, so a monolingual model would forfeit half the data. Second,
predictions are scored by mean absolute error (MAE) on a five-point scale,
making the task \emph{ordinal}---confusing a one-star review with a
five-star one is four times as costly as confusing it with a two-star one,
unlike standard classification where every error is equal. Third, the text
is \emph{noisy}: it contains typos, emojis, repeated characters
(\textit{sooooo}), sarcasm, and genuinely ambiguous reviews.

A natural question is which representation of the text best captures
sentiment under these constraints. We answer it empirically by comparing
three families of increasing richness---sparse lexical features, frozen
contextual embeddings, and task-adapted (fine-tuned) multilingual
transformers---and find that adapting the encoder to the task is the
decisive factor. We then show that two cheap, ordinal-aware techniques
push performance further: a decoding rule matched to the MAE metric, and an
ensemble of \emph{diverse} encoders that outperforms any single base model
and sidesteps the instability of training a larger one on commodity
hardware.

Our contributions are:
\begin{itemize}
    \item A controlled comparison of sparse lexical features, frozen
    contextual embeddings, and a fine-tuned multilingual encoder under an
    identical data split and evaluation protocol
    (\cref{sec:baselines,sec:main}).
    \item An MAE-optimal decoding rule that replaces $\arg\max$ with the
    median of the predictive distribution, lowering MAE at no training cost
    (\cref{sec:decoding}).
    \item A diversity ensemble of two multilingual encoders that exceeds the
    strong single-model baseline (\cref{sec:ensemble}), motivated by the
    observation that scaling to a single larger model is numerically
    unstable in half precision on our GPUs (\cref{sec:discussion}).
\end{itemize}

\section{Models and Methods}
\label{sec:methods}

\subsection{Data and exploratory analysis}
\label{sec:data}
The training set contains $252{,}000$ labeled reviews and the test set
$168{,}000$. The five classes are \emph{exactly balanced} ($50{,}400$ each),
so we use an unweighted loss and report accuracy alongside MAE without
class re-weighting. All model selection uses a single $10\%$ stratified
hold-out split with a fixed seed; we never tune on the public leaderboard,
following standard practice to avoid overfitting to it.

An exploratory pass informs three design choices. (i)~A language heuristic
estimates $\approx\!52\%$ English and $\approx\!48\%$ German, with the split
uniform across labels; we therefore use a \emph{multilingual} encoder
rather than translating or training two monolingual models. (ii)~Each
review is a short title followed by a body; the $99$th percentile length is
$161$ words, so a maximum sequence length of $256$ tokens covers virtually
all reviews while bounding compute. (iii)~Surface cues flagged as
informative behave non-trivially: the frequency of ``\texttt{!}'' is
\emph{U-shaped} in the rating (high for both one- and five-star reviews),
so it signals \emph{intensity} rather than polarity. A model that relies on
such cues alone would conflate the two extremes---the worst possible error
under MAE---motivating representations that capture context.

\subsection{Evaluation metric}
\label{sec:metric}
Ratings are encoded as $y \in \{0,\dots,4\}$ and scored by
\begin{equation}
\label{eq:score}
    L(\hat{\mathbf y}, \mathbf y) = 1 - \frac{1}{4n}\sum_{i=1}^{n} |y_i - \hat y_i|
    = 1 - \frac{\mathrm{MAE}}{4},
\end{equation}
so the score lies in $[0,1]$ and larger ordinal errors are penalized
proportionally. This metric drives both our loss-agnostic decoding rule
(\cref{sec:decoding}) and our emphasis on avoiding extreme misclassifications.

\subsection{Baselines}
\label{sec:baselines}
\textbf{B1: TF--IDF + logistic regression.} We vectorize each review with
TF--IDF over unigrams and bigrams ($50{,}000$ features, sublinear term
frequency) and fit a multinomial logistic regression. This captures
predictive lexical tokens directly but ignores word order and context.

\textbf{B2: Frozen encoder + logistic regression.} We encode each review
with a frozen multilingual sentence transformer
(\texttt{paraphrase-multilingual-MiniLM-L12-v2}, $384$ dimensions) and fit
the same logistic regression on the embeddings. This \emph{linear probe}
isolates the value of pre-trained contextual semantics \emph{without} task
adaptation, and contrasts directly with B1 (dense vs.\ sparse) and with our
main method (frozen vs.\ fine-tuned).

\subsection{Fine-tuned multilingual transformers}
\label{sec:main}
Our main models fine-tune a pre-trained multilingual encoder end-to-end
with a five-way classification head. We use two architecturally distinct
encoders of comparable size: XLM-RoBERTa-base~\citep{conneau2020xlmr} and
mDeBERTa-v3-base, the latter using disentangled attention. Both are trained
with AdamW (learning rate $2\times10^{-5}$, weight decay $0.01$, linear
warmup over $6\%$ of steps), a maximum length of $256$ tokens, batch size
$16$, for three epochs in mixed precision, selecting the checkpoint with the
best validation score. Gradients are clipped to unit norm.

\subsection{MAE-optimal decoding}
\label{sec:decoding}
Given predicted class probabilities $p = (p_0,\dots,p_4)$, the usual
$\arg\max$ returns the \emph{mode}, which is optimal for $0/1$ loss. Under
absolute-error loss the optimal point estimate is instead the \emph{median}
of the predictive distribution---the smallest $k$ with
$\sum_{j\le k} p_j \ge 0.5$. Because our metric (\cref{eq:score}) is
MAE-based, we decode with the median for \emph{every} model. The rule needs
no retraining and consistently reduces MAE: it trades a slightly lower exact
accuracy for fewer large, costly errors.

\subsection{Diversity ensemble}
\label{sec:ensemble}
Scaling to a single larger encoder (XLM-RoBERTa-large) is attractive but, on
our hardware, half-precision training diverges and full precision does not
fit in memory (\cref{sec:discussion}). Instead we exploit
\emph{architectural diversity}: we average the predictive distributions of
the two fine-tuned base encoders and apply the median decoder of
\cref{sec:decoding}. Averaging probabilities of models that make
\emph{different} errors reduces variance and improves the score over either
model alone, at the cost of one additional fine-tuning run rather than a
larger, unstable model.

\section{Results}
\label{sec:results}
% TODO: fill table with final numbers. Val from local split; LB from Kaggle.

\begin{table}[t]
\caption{Validation and public-leaderboard scores ($1-\mathrm{MAE}/4$).}
\label{tab:main}
\vskip 0.1in
\begin{center}
\begin{small}
\begin{tabular}{lcc}
\toprule
Method & Val score & Public LB \\
\midrule
B1: TF--IDF + LogReg              & 0.882 & TODO \\
B2: Frozen enc.\ + LogReg         & 0.850 & TODO \\
XLM-R-base (fine-tuned)           & 0.905 & TODO \\
mDeBERTa-v3-base (fine-tuned)     & TODO  & TODO \\
\textbf{Ensemble (ours)}          & \textbf{TODO} & \textbf{TODO} \\
\bottomrule
\end{tabular}
\end{small}
\end{center}
\vskip -0.1in
\end{table}

% TODO: ablation table (argmax vs median; base vs large), per-language MAE,
% confusion matrix figure.

\section{Discussion}
\label{sec:discussion}
% TODO
% - Why B2 < B1: frozen paraphrase embeddings encode general semantics, not
%   sentiment-specific lexical signal; TF-IDF captures predictive tokens
%   directly. Motivates fine-tuning.
% - Error analysis: residuals concentrate on sarcasm / ambiguous reviews.
% - Cost/benefit of large vs base (compute vs score).

\section{Summary}
\label{sec:summary}
% TODO: 3-4 sentences restating contributions in light of results.

\section*{Declaration of Originality}
% Required. List pre-trained models, libraries, and AI assistance used.
We confirm this work is our own. We used the publicly available pre-trained
models \texttt{xlm-roberta-*} and
\texttt{paraphrase-multilingual-MiniLM-L12-v2}, and the libraries
\texttt{scikit-learn}, \texttt{PyTorch}, and \texttt{transformers}.
% TODO: add AI-assistance disclosure per course policy.

\bibliography{report}
\bibliographystyle{icml2024}

\end{document}
